La constante de Catalan suele aparecer como una serie alternada:\par

\[
G
=
1-\frac1{3^2}+\frac1{5^2}-\frac1{7^2}+\cdots.
\]

Esa fórmula da su valor y la genealogía HMT explica por qué existe esa alternancia concreta.\par

En HMT, el origen está en el cuarto de torsión.\par

Existe un operador interno \(J\) tal que\par

\[
J^2=-I.
\]

La raíz de \(-1\) emerge de dos aplicaciones sucesivas de un cuarto de giro. La estructura compleja es la memoria algebraica de una rotación interna.\par

Ese giro genera un carácter de orden cuatro. Cuando se publica sobre los enteros positivos, produce exactamente el patrón residual que alterna entre \(+1\), \(0\), \(-1\), \(0\).\par

La constante de Catalan aparece entonces como traza espectral:\par

\[
G
=
L(2,\chi_{-4})
=
\operatorname{Tr}(N^{-2}Q_4).
\]

Esto cambia su significado. Catalan es simultáneamente la suma de una serie particular y la respuesta espectral de los enteros ante un cuarto de torsión.\par

Cada término de la serie registra cómo el giro interno actúa sobre una clase residual. La alternancia es la sombra aritmética de la rotación.\par

La estructura triádica posee su propio carácter módulo \(3\). Al combinarlo con el carácter de orden cuatro, el teorema chino del resto produce un carácter módulo \(12\):\par

\[
\chi_{12}
=
\chi_{-3}\chi_{-4}.
\]

El módulo \(12\) aparece como el espacio común mínimo en el que pueden coexistir la orientación triádica y el cuarto de torsión.\par

La dodecafase es, en ese sentido, una reconciliación de \(3\) y \(4\).\par

Del carácter combinado emergen\par

\[
L(1,\chi_{12})
=
\frac{\log(2+\sqrt3)}{\sqrt3},
\]

y\par

\[
L(2,\chi_{12})
=
\frac{\pi^2}{6\sqrt3}.
\]

La cantidad\par

\[
2+\sqrt3
\]

es la unidad fundamental del campo real \(\mathbb Q(\sqrt3)\). Su logaritmo representa un desplazamiento hiperbólico. De nuevo, una estructura finita de residuos abre una dinámica multiplicativa infinita.\par

Aquí se encuentran giro, torsión, carácter aritmético, unidad de Pell y desplazamiento hiperbólico.\par

La constante de Euler–Kronecker del campo \(\mathbb Q(\sqrt3)\) prolonga esta estructura. Mide el término regular que permanece tras separar el polo principal de la función zeta del campo. En la genealogía HMT, este término constituye el invariante regular de memoria aritmética fina que subsiste después de retirar la divergencia universal.\par

La familia Bendersky–Glaisher prolonga el mecanismo en una torre. Las derivadas espectrales de zeta generan esa torre. Apéry y otras constantes asociadas a valores impares se organizan como miembros de una misma familia funcional.\par

La aportación novedosa reside en la genealogía que reúne las constantes clásicas:\par

\begin{itemize}[leftmargin=*,itemsep=0.35em]
\item el cuarto de torsión produce la estructura compleja;
\item el residuo triádico produce la estructura módulo \(3\);
\item su composición produce el carácter módulo \(12\);
\item el carácter genera valores \(L\);
\item los valores \(L\) abren el campo cuadrático y sus constantes espectrales.
\end{itemize}

La torsión geométrica se ha convertido en aritmética conservando íntegramente su orientación.\par
